%==============================================================
%  frontmatter/preface.tex — Avant-propos
%==============================================================

\chapter*{Avant-propos}
\addcontentsline{toc}{chapter}{Avant-propos}
\thispagestyle{plain}

Ce recueil rassemble des leçons sur les théorèmes incontournables des
mathématiques de classes préparatoires. L'ambition
n'est pas l'exhaustivité : il s'agit de mettre au propre, avec rigueur et
clarté, des résultats fondamentaux :  énoncé, intuition, démonstration 
complète, exemples et exercices.

\medskip
\noindent\textbf{Structure d'une leçon.} Chaque fiche suit le même plan :
\begin{enumerate}
  \item \textbf{Énoncé} — le théorème dans toute sa précision, encadré en
        bleu avec les hypothèses explicitées.
  \item \textbf{Signification intuitive} — une explication qualitative en
        quelques lignes, sans formalisme.
  \item \textbf{Idée de la preuve} — la stratégie en prose, avant tout calcul.
  \item \textbf{Démonstration formelle} — la preuve complète, structurée
        en étapes numérotées, avec une barre latérale bleue.
  \item \textbf{Exemples \& contre-exemples} — pour ancrer le résultat et
        cerner ses limites.
  \item \textbf{Exercices} — des problèmes classiques pour s'entraîner.
\end{enumerate}

\medskip
\noindent\textbf{Organisation.} Les leçons sont regroupées en six grandes parties
thématiques : Analyse réelle, Structures algébriques, Arithmétique, Combinatoire,
Topologie et Algèbre linéaire. Cette structure reflète celle des programmes de
classes préparatoires scientifiques.

\medskip
\noindent\textbf{Comment utiliser ce recueil.} Chaque leçon est autonome
et peut être lue indépendamment. L'index en fin de volume permet de retrouver
rapidement un concept ou un théorème. Le glossaire rappelle les définitions
essentielles. Les références bibliographiques orientent vers les sources
classiques pour approfondir.

\vspace{1cm}
\begin{flushright}
  \textit{Bonne lecture et bon courage.}
\end{flushright}

\clearpage
